\section{Conclusions}

A count of local government areas (LGAs) above a prevalence threshold is a distinct posterior quantity, not a list of area estimates. Its variance depends on pairwise exceedance probabilities, which marginal-only reporting does not supply. At the 50\% threshold, the fitted expected count was 115 of 774 LGAs, with a 95\% posterior interval of 101--130. The pairwise term means that an independence calculation would understate the fitted count's standard deviation by 14\%. This interval depends on the fitted joint dependence structure and was not validated by the area-level cross-validation.

The marginal probabilities allow the count's standard deviation to range from 0.156 to 74.11. The model's variance sits only 0.99\% of the way through the feasible variance range. The posterior interval for the sample-analogue identified count runs from 87 to 231. Of its 144-area width, 120 reflects unrestricted contributions from LGAs without sampled clusters. Those LGAs contain 11.04\% of the population, and the lower endpoint for sampled areas still uses model-smoothed prevalence estimates.

Spatial cross-validation used pass rules fixed before fitting. At the exposure faced by unclustered LGAs, area-level coverage was 0.890 against the nominal 0.95; withholding a whole region yielded 0.775. An interval-widening factor \(k=1.178\) was checked out of fold for area-level intervals. The protocol is reusable because it replaces an assumed nominal level with measured performance in prediction settings that match intended use.

The transferable parts are the sharp bounds and the pre-declared protocol, not a new estimator. The Nigerian map is a demonstration, and its area-level precision is supported only for prediction settings represented by sampled data.
